\begin{lema}[Prescripción de índices y llegada larga]\label{lem:prescripcion-larga}
Sea \(S_M\) la entrada histórica regular en \(M=\lceil\delta n\rceil\) y \(H^0_j\)
la supervivencia regular hasta \(j\) inclusive después de esa entrada. Si se prescriben
\(s\) filas y \(t\) columnas, con \(s,t\le j/2\), su probabilidad de seguir vivas junto
con \(H_j^0\), dada la historia en \(M\) sobre \(S_M\), es a lo sumo
\[
 (j/M)^{s+t}\exp(C_{\rm pres}(st+t)/j).
\]
Bajo la condición sobre \(x_0\) del Lema~\ref{lem:llegada-integrada},
\(\mathbb E_\nu[1+X;H^0_j\setminus H_j]=O(e^{-2j})\), con \(X=Qe^{-L}\).
\end{lema}
\begin{proof}
Si falta una marca en \(M\), la probabilidad es cero. En cada padre regular de tamaño
\(k>j\), la fila uniforme toca alguna de las \(s\) marcas con probabilidad \(s/k\).
Dada la fila, la transición vale \(p_{il}(1+O(1/k))\) por la prueba del
Teorema~8.1 de~\cite{pc}, con constantes de clase independientes de \(n/k\).
Las sumas de columna de \(p\) valen uno. Así, la probabilidad de una de las \(t\) columnas
es \(\ge t/k-Ct/k^2\), y la de tocar ambos conjuntos es \(\le Cst/k^2\) por difusión.
Evitar ambas marcas y seguir regular tiene probabilidad a lo sumo
\((1-s/k)(1-t/k)\exp(C(st+t)/k^2)\), pues el producto es al menos \(1/4\).
Al iterar, \(\prod_{k=j+1}^M(1-s/k)=(j)_s/(M)_s\le(j/M)^s\);
las columnas dan el otro factor y \(\sum_{k>j}k^{-2}\le1/j\).

En una diagonal hay a lo sumo \(M\) candidatos a la entrada. Si quedan \(L\) sitios,
para \(r=\lfloor\alpha j/4\rfloor\) y \(R=\lceil\alpha j\rceil\) el indicador de \(L\ge R\)
está dominado por \((L)_r/(R)_r\). Cada lista ordenada exige \(r\) filas y \(r\) columnas
distintas. Como \((R)_r\ge(\alpha j/2)^r\), al sumar en menos de \(4n\) diagonales resulta
\[
 \nu(H^0_j\setminus H_j)\le
 4n\left(\frac{2j}{\alpha M}\right)^r
      \exp(C_{\rm pres}(r^2+r)/j).
\]
Al dividir su logaritmo por \(j\), el límite superior es a lo sumo
\((\alpha/4)\log x_0+C_{\rm pres}\alpha^2/16<-(\beta+3)\), pues
\(j=\lfloor bn\rfloor\) y \(b\) es fijo. Desde \(n_0(b)\) la masa es
\(\le e^{-(\beta+2)j}\). En las llegadas regulares, las cotas estáticas de
\cite[6.5--6.8]{tc} dan \(X\le e^{C_0+\beta j}\); para un tablero incompletable \(X=0\).
Multiplicar por \(1+X\) da el coste indicado. El margen \(3\) del exponente cubre
el redondeo y el factor de unión antes de pagar \(\beta j\).
\end{proof}
